% !TeX root = ../main.tex
% ============================================================
%  sections/08_conclusion.tex
% ============================================================

\section{Conclusion}
\label{sec:conclusion}

This paper has addressed the design of sustainable closed-loop
supply chains under a carbon-tax policy by modelling the interaction
between the regulator and the supply-chain operator as a bilevel
mixed-integer programme. The regulator lexicographically minimises
emissions and then the tax rate, while the operator solves a
multi-echelon, multi-period \CLSC design and operations problem with
integer facility, technology, capacity, and transportation decisions.
We developed a deterministic bounding procedure (\DBP) that combines
a relaxation-based lower bound with exact follower evaluations and
returns validated bilevel solutions on the generated instances.

\subsection{Summary of findings}
\label{sec:conc:summary}

Three findings emerge from the computational study:
\begin{enumerate}
  \item \textbf{The emissions-minimising tax is instance-specific and
        sits on a saturating frontier.} Optimal rates span
        $20$--$200$ USD/t\,CO\textsubscript{2}e, and beyond a
        moderate level extra tax yields negligible abatement, so a
        purely emissions-driven regulator is ill-posed unless paired
        with a cost or meaningful-abatement criterion.
  \item \textbf{Technology, not transport, is the binding abatement
        lever.} Expanding the production and remanufacturing
        technology menu cuts optimal emissions by $26$--$33\%$, while
        expanding the transport-mode catalogue is essentially inert;
        emissions are driven by demand and by production and upstream
        (scope-3) emission intensities.
  \item \textbf{The \DBP is reliable but horizon-bound.} All 95
        committed runs are validated at relative gaps no larger than
        $1.9\times10^{-5}$ in two to three iterations; T3--T6
        instances solve within hours, whereas the T12 instances are
        intractable within the compute budget, marking horizon length
        as the main computational bottleneck.
\end{enumerate}

\subsection{Academic contributions}
\label{sec:conc:contributions}

The contributions of the paper are threefold. First, it formulates a
bilevel \CLSC model that integrates a discrete, lexicographic
regulatory tax decision with a detailed integer supply-chain design
problem. Second, it develops a convergent \DBP with saved-solution
validation and uses it to derive policy- and management-relevant
insight --- the saturating emission--tax frontier and the primacy of
the technology lever --- that single-level models cannot reveal.
Third, it provides a tested generated benchmark, a sensitivity and
scenario suite, and a documented result schema that ties every
reported number to saved solver outputs.

\subsection{Limitations and future research directions}
\label{sec:conc:future}

Several extensions deserve attention.

\paragraph{Cost-aware leader objectives.}
The saturating frontier makes a purely emissions-minimising leader
ill-posed; a natural next step is a leader objective that trades
emissions against tax burden explicitly, or that imposes a minimum
meaningful-abatement threshold, so the optimal tax is well-identified
without relying on the lexicographic tie-break.

\paragraph{Scalability to longer horizons.}
The T12 instances are out of reach within the present compute budget;
decomposition, warm starts, and parallel \DBP implementations are the
most promising routes to the longest horizons and larger networks.

\paragraph{External benchmark implementations.}
KKT reformulations, exhaustive bilevel grid search, or other
benchmark methods should be implemented before any cross-method
performance claims are made.

\paragraph{Mechanism-level diagnostics.}
The scenario design isolates the technology lever but not the
individual mechanisms within it; decomposing abatement by facility,
technology, remanufacturing, inventory, and transport decision would
support finer mechanism-specific claims.

\paragraph{Stochastic and robust extensions.}
Embedding demand, return-rate, cost, and emission uncertainty into
the follower problem would broaden the model's applicability to
volatile markets.

\paragraph{Pessimistic and risk-averse leader models.}
Implementing pessimistic tie-breaking and risk-averse leader
objectives would strengthen the regulatory interpretation when the regulator
cannot assume favourable follower tie-breaking.

\paragraph{Mixed regulatory instruments.}
Caps, subsidies, cap-and-trade, and dynamic tax trajectories are
natural extensions of the leader programme.

\paragraph{Algorithmic enhancements.}
Adaptive sampling, decomposition, warm starts, and parallel \DBP
implementations are promising routes for larger instances.

Pursuing these extensions would turn the current reproducible
prototype into a broader decision-support framework for joint supply
chain and climate-policy design.

\section*{Acknowledgements}
The authors thank the anonymous reviewers for their constructive
comments. This work was supported by [Funding Agency, Grant Number].

\section*{Author Contributions}
\textbf{Marcelo Pinto}: Conceptualisation, Methodology, Software,
Formal Analysis, Writing -- Original Draft.
\textbf{Paulo Nascimento}: Methodology, Validation, Writing --
Review \& Editing.
\textbf{Crist\'ov\~ao Silva}: Supervision, Writing -- Review \&
Editing.
\textbf{Maria Jo\~ao Alves}: Methodology, Formal Analysis, Writing
-- Review \& Editing.
\textbf{Carlos Henggeler}: Supervision, Writing -- Review \&
Editing.
\textbf{Samuel Moniz}: Conceptualisation, Supervision, Project
Administration, Funding Acquisition, Writing -- Review \& Editing.

\section*{Declaration of Competing Interests}
The authors declare that they have no known competing financial
interests or personal relationships that could have appeared to
influence the work reported in this paper.
